% sec:adm-extrinsic-curvature — Extrinsic Curvature
\section{Extrinsic Curvature}
\label{sec:adm-extrinsic-curvature}

The spatial metric $\spatialmetric$ captures the geometry within each
slice $\Sigma_t$ --- distances, angles, and intrinsic curvature --- but
says nothing about how the slice sits inside the surrounding spacetime.
A flat sheet of paper and a cylinder share the same intrinsic geometry
(both have vanishing Gaussian curvature), yet one is flat and the other
curves: their \emph{embeddings} differ.  The analogy is imperfect: a
cylinder's extrinsic curvature is a scalar (a single principal
curvature), whereas $\extrinsic$ is a six-component tensor that encodes
direction-dependent bending --- the slice can simultaneously expand in
one direction and contract in another.  The extrinsic curvature
quantifies this distinction for the spatial slices of the 3+1
decomposition.  Together with $\spatialmetric$, it constitutes the
complete gravitational initial data whose evolution
\cref{sec:adm-equations} will govern.
\coderef{ADMSpacetime}{extrinsicCurvature}

\paragraph{Definition.}

Two neighbouring Eulerian observers separated by a displacement
$\delta x^i$ on $\Sigma_t$ both follow the unit normal $n^\mu$ to the
next slice.  Because the normals at nearby points need not be parallel,
the observers converge or diverge as they advance --- the extrinsic
curvature measures this rate of convergence.  Defining it as half the Lie
derivative of the induced metric along the normal,
%
\begin{equation}
  K_{\mu\nu} = -\frac{1}{2}\,\lie{n}\gamma_{\mu\nu} \,,
  \label{eq:adm-extrinsic-def}
\end{equation}
%
makes the geometric content transparent: if the extrinsic curvature is
positive in some direction, two Eulerian observers separated in that
direction see their proper distance shrinking as they advance to the
next slice --- the slice is locally concave, bending toward the future.
The minus sign is conventional --- positive trace $\meanextrinsic$
corresponds to overall contraction.  Since $\gamma_{\mu\nu}$ is
symmetric and the Lie derivative preserves tensor symmetry,
$K_{\mu\nu} = K_{\nu\mu}$ by construction.
\physnote{The extrinsic curvature is the gravitational analogue of the
second fundamental form in classical differential geometry --- it measures
how the embedded surface bends within the ambient space.}

\paragraph{Covariant-derivative form.}

Expanding the Lie derivative in \cref{eq:adm-extrinsic-def} yields a
formula better suited to derivations.  Using
$\gamma_{\mu\nu} = \metric + n_\mu\,n_\nu$
(\cref{eq:adm-induced-metric}) and metric compatibility
$\cd{\alpha}\metric = 0$, the Lie derivative of the spacetime metric
reduces to $\lie{n}\metric = \cd{\mu}n_\nu + \cd{\nu}n_\mu$.  The
remaining terms from $\lie{n}(n_\mu\,n_\nu)$ involve the acceleration
$a_\mu = n^\alpha\,\cd{\alpha}n_\mu$ of the Eulerian observer, giving
the full expansion
%
\begin{equation}
  \lie{n}\gamma_{\mu\nu}
  = \cd{\mu}n_\nu + \cd{\nu}n_\mu
    + a_\mu\,n_\nu + n_\mu\,a_\nu \,.
  \label{eq:adm-lie-gamma-expand}
\end{equation}
%
We now need the spatial covariant derivative $D_\mu$ --- the intrinsic
derivative operator on $\Sigma_t$ that will appear throughout the rest
of the 3+1 formalism.  It is the unique torsion-free connection
compatible with $\spatialmetric$, built from the three-dimensional
Christoffel symbols $\spatialchristoffel{i}{j}{k}$.  For a spatial
scalar $f$, the spatial gradient coincides with the full gradient
projected onto $\Sigma_t$:
$D_\mu f = \gamma^\nu{}_\mu\,\cd{\nu}f$.

Using this projection, the acceleration reduces to a pure spatial
gradient.  Differentiating
$n_\mu = -\lapse\,\cd{\mu}t$
(\cref{eq:adm-normal-covector}) gives
$\cd{\alpha}n_\mu
= -(\cd{\alpha}\lapse)\,\cd{\mu}t
  - \lapse\,\cd{\alpha}\cd{\mu}t$.
Contracting with $n^\alpha$ and using
$n^\alpha\,\cd{\alpha}t = -1/\lapse$ (from the normalisation
$n_\mu = -\lapse\,\cd{\mu}t$ contracted with $n^\mu$), the
$\cd{\alpha}\cd{\mu}t$ term produces $\cd{\mu}\lapse/\lapse$ after
symmetry of second derivatives and projection, yielding
%
\begin{equation}
  a_\mu = D_\mu \ln\lapse \,.
  \label{eq:adm-acceleration}
\end{equation}
%
The acceleration is spatial
($a_\mu\,n^\mu = 0$, from differentiating $n_\mu\,n^\mu = -1$), so
the terms $a_\mu\,n_\nu + n_\mu\,a_\nu$ in
\cref{eq:adm-lie-gamma-expand} contribute only mixed normal--tangent
components.

To extract the spatial content, project both indices of
\cref{eq:adm-extrinsic-def} onto $\Sigma_t$.  The double projection of
\cref{eq:adm-lie-gamma-expand} kills every term containing an $n_\mu$
or $n_\nu$ factor, since
$\gamma^\alpha{}_\mu\,n_\mu = 0$
(\cref{eq:adm-projector-orthogonal}), leaving
$K_{\mu\nu}
= -\gamma^\alpha{}_\mu\,\gamma^\beta{}_\nu\,\cd{\alpha}n_\beta$.
But the second projector is redundant: differentiating the
normalisation $n_\beta\,n^\beta = -1$ gives
$n^\beta\,\cd{\alpha}n_\beta = 0$, so $\cd{\alpha}n_\beta$ already has
no component along $n^\beta$, and
$\gamma^\beta{}_\nu\,\cd{\alpha}n_\beta = \cd{\alpha}n_\nu$.
The result is
%
\begin{equation}
  K_{\mu\nu} = -\gamma^\alpha{}_\mu\,\cd{\alpha}n_\nu \,.
  \label{eq:adm-extrinsic-covariant}
\end{equation}
%
The projector $\gamma^\alpha{}_\mu$ strips out the normal component of
$\cd{\alpha}n_\nu$, leaving only the part that describes how the
normal rotates within the slice.  Note that
$\gamma^\alpha{}_\mu\,\cd{\alpha}n_\nu$ is symmetric in $\mu,\nu$
despite appearances: $n_\mu$ is hypersurface-orthogonal, so
$\cd{[\mu}n_{\nu]} = -n_{[\mu}\,a_{\nu]}$, and the antisymmetric part
of $\cd{\mu}n_\nu$ is purely normal, annihilated by the projector.
\warnnote{The sign convention
$K_{\mu\nu} = -\gamma^\alpha{}_\mu\,\cd{\alpha}n_\nu$ (with the minus
sign) matches Baumgarte \& Shapiro \cite{Baumgarte2010} and Alcubierre
\cite{Alcubierre2008}.  Wald \cite{Wald1984} uses the opposite
convention --- always check signs when comparing across sources.}

\paragraph{Purely spatial.}

For $\extrinsic$ to serve as initial data on a spatial slice, it must
carry no information about the normal direction --- it must be purely
spatial.  Contracting \cref{eq:adm-extrinsic-covariant} with $n^\mu$
gives zero from the projector: $K_{\mu\nu}\,n^\mu
= -\gamma^\alpha{}_\mu\,n^\mu\,\cd{\alpha}n_\nu = 0$, since
$\gamma^\alpha{}_\mu\,n^\mu = 0$
(\cref{eq:adm-projector-orthogonal}).  Contracting instead with $n^\nu$
also vanishes: as shown above,
$n^\nu\,\cd{\alpha}n_\nu = 0$, so
$K_{\mu\nu}\,n^\nu
= -\gamma^\alpha{}_\mu\,n^\nu\,\cd{\alpha}n_\nu = 0$.  The extrinsic
curvature therefore has only six independent components --- the spatial
components $\extrinsic$.  These six components, together with the six of
$\spatialmetric$, constitute the twelve pieces of gravitational initial
data that the constraint and evolution equations of
\cref{sec:adm-equations} will govern.

\begin{figure}[htbp]
  \centering
  % adm-extrinsic-curvature.tex — Convergence/divergence of normals
% Inputable TikZ fragment for fig:adm-extrinsic-curvature
% Requires: tikz with calc, arrows.meta (loaded by ehbook.cls)
%
% Cross-section of a curved spatial slice Sigma_t with unit normals
% at several points.  Left trough (concave upward): normals converge
% toward the future, K > 0.  Right peak (concave downward): normals
% diverge, K < 0.
%
% Normal vectors are hand-computed perpendiculars to the curve.
% At the trough (x≈2.5) the tangent is horizontal, so the normal
% is vertical; on either side the tangent tilts, rotating the normal
% toward the centre of curvature (convergence) or away from it
% (divergence).

\begin{tikzpicture}[
    >={Stealth[length=5pt]},
    font=\footnotesize,
  ]

  % ── Future direction indicator ──────────────────────────────────
  \draw[->, EHMarginGray, thin]
    (-0.2, 0.8) -- (-0.2, 3.2);
  \node[font=\scriptsize, text=EHMarginGray, rotate=90, anchor=south]
    at (-0.45, 2.0) {future};


  % ── Curved spatial slice Σ_t ────────────────────────────────────
  %
  % S-shaped cross section:
  %   Trough at x ≈ 2.5  →  concave up  →  K > 0  (normals converge)
  %   Peak   at x ≈ 5.5  →  concave down →  K < 0  (normals diverge)

  \draw[EHMarginGray, semithick]
    plot[smooth, tension=0.6] coordinates {
      (0.5, 2.50) (1.5, 1.70) (2.5, 1.20)
      (3.5, 1.80) (4.0, 2.50)
      (4.5, 3.20) (5.5, 3.80) (6.5, 3.20) (7.5, 2.50)};

  \node[font=\small, text=EHMarginGray, anchor=west]
    at (7.70, 2.50) {$\Sigma_t$};


  % ── Concave region normals (K > 0, converging) ─────────────────
  %
  % Three normals tilt inward toward the centre of curvature
  % (above the trough), so their tips are closer together than
  % their bases.

  % Normal at x = 1.5 — tilts right
  \coordinate (N1) at (1.50, 1.70);
  \draw[->, EHJetBlue, very thick, shorten >=1pt]
    (N1) -- ++(0.38, 0.85);

  % Normal at x = 2.5 (trough) — points straight up
  \coordinate (N2) at (2.50, 1.20);
  \draw[->, EHJetBlue, very thick, shorten >=1pt]
    (N2) -- ++(0.00, 0.93);

  % Normal at x = 3.5 — tilts left
  \coordinate (N3) at (3.50, 1.80);
  \draw[->, EHJetBlue, very thick, shorten >=1pt]
    (N3) -- ++(-0.35, 0.86);

  % Label one normal
  \node[font=\scriptsize, text=EHJetBlue, anchor=south west, inner sep=2pt]
    at (2.55, 1.88) {$n^\mu$};


  % ── Convex region normals (K < 0, diverging) ───────────────────
  %
  % Three normals tilt outward away from the centre of curvature
  % (below the peak), so their tips are further apart than their
  % bases.

  % Normal at x = 4.5 — tilts left
  \coordinate (N4) at (4.50, 3.20);
  \draw[->, EHJetBlue, very thick, shorten >=1pt]
    (N4) -- ++(-0.38, 0.85);

  % Normal at x = 5.5 (peak) — points straight up
  \coordinate (N5) at (5.50, 3.80);
  \draw[->, EHJetBlue, very thick, shorten >=1pt]
    (N5) -- ++(0.00, 0.93);

  % Normal at x = 6.5 — tilts right
  \coordinate (N6) at (6.50, 3.20);
  \draw[->, EHJetBlue, very thick, shorten >=1pt]
    (N6) -- ++(0.38, 0.85);


  % ── Base-point markers (drawn last, on top) ────────────────────
  \foreach \n in {N1,N2,N3,N4,N5,N6}
    \fill (\n) circle (1.2pt);


  % ── Region annotations ─────────────────────────────────────────
  %
  % K > 0 label below the trough (in the concave bowl)
  \node[font=\small, text=EHAccretionGold]
    at (2.50, 0.55) {$K > 0$};
  \node[font=\scriptsize, text=EHMarginGray]
    at (2.50, 0.18) {normals converge};

  % K < 0 label above the peak (where normals fan out)
  \node[font=\small, text=EHAccretionGold]
    at (5.50, 5.05) {$K < 0$};
  \node[font=\scriptsize, text=EHMarginGray]
    at (5.50, 5.40) {normals diverge};

\end{tikzpicture}

  \caption[Convergence of normals and extrinsic curvature]{%
    Extrinsic curvature as the convergence of normals to a spatial
    slice.  On the concave portion of $\Sigma_t$ (left), the unit
    normals $n^\mu$ drawn at neighbouring points converge toward the
    future, corresponding to positive mean curvature $K > 0$.  On the
    convex portion (right), the normals diverge, giving $K < 0$.  The
    individual components of $K_{ij}$ carry the directional
    information; the trace $K$ records only the net volume change
    along $n^\mu$.}
  \label{fig:adm-extrinsic-curvature}
\end{figure}

\paragraph{Relation to the time derivative of the spatial metric.}

The Lie-derivative definition \cref{eq:adm-extrinsic-def} links
$\extrinsic$ to the rate of change of spatial distances as measured by
Eulerian observers, but a numerical code works with the coordinate time
derivative $\partial_t\gamma_{ij}$.  Converting between the two requires
the decomposition of the unit normal in adapted coordinates.  From
$(\partial/\partial t)^\mu = \lapse\,n^\mu + \beta^\mu$
(\cref{eq:adm-time-vector}), solving for $n^\mu$ gives
%
\begin{equation}
  n^\mu = \frac{1}{\lapse}\bigl(1,\,-\shift\bigr) \,,
  \label{eq:adm-normal-adapted}
\end{equation}
%
so the Lie derivative along $n$ decomposes as
$\lie{n} = (1/\lapse)(\lie{\partial_t} - \lie{\beta})$.  In coordinates
adapted to $\partial_t$, the Lie derivative
$\lie{\partial_t}\gamma_{ij} = \partial_t\gamma_{ij}$.  The Lie
derivative of a metric along a vector field satisfies the standard
identity $\lie{\beta}\gamma_{ij} = D_i\beta_j + D_j\beta_i$, where
$\beta_i = \gamma_{ij}\,\beta^j$ is the shift with its index lowered by
the spatial metric.  Substituting into \cref{eq:adm-extrinsic-def}
yields the evolution equation for the spatial metric:
%
\begin{equation}
  \partial_t \gamma_{ij}
  = -2\,\lapse\,\extrinsic
    + D_i\beta_j + D_j\beta_i \,.
  \label{eq:adm-metric-evolution}
\end{equation}
%
This is one half of the ADM evolution system.  It tells us how the
spatial geometry changes from one slice to the next: the first term
records the genuine stretching due to the extrinsic curvature
(weighted by the lapse), while the shift terms account for the
coordinate drift of the spatial grid.  Setting $\shift = 0$ and
$\lapse = 1$ (geodesic slicing) reduces this to
$\partial_t\gamma_{ij} = -2\,\extrinsic$ --- the extrinsic curvature is
half the coordinate-time rate of change of the spatial metric.  The
companion equation governing $\partial_t \extrinsic$ follows from the
Gauss--Codazzi decomposition and is derived in \cref{sec:adm-equations}.
\implnote{In the codebase, \cref{eq:adm-metric-evolution} is used
both to initialise $\extrinsic$ from an analytic time derivative of
$\spatialmetric$ and to advance $\spatialmetric$ during evolution.}

\paragraph{Mean curvature.}

In many situations --- singularity avoidance, constraint equations,
gauge choices --- a single scalar summary of whether the slice is
expanding or contracting matters more than the full tensor.  Tracing the
extrinsic curvature with the inverse spatial metric defines the mean
curvature
%
\begin{equation}
  \meanextrinsic = \inversespatialmetric\,\extrinsic \,,
  \label{eq:adm-mean-curvature}
\end{equation}
%
a scalar that measures the fractional rate of change of the spatial
volume element along the normal.  Tracing the definition
\cref{eq:adm-extrinsic-def} with $\inversespatialmetric$ and using the
Jacobi identity for the Lie derivative of a determinant
($\gamma^{ij}\,\lie{n}\gamma_{ij} = \lie{n}\ln\gamma$) gives
$\meanextrinsic = -\lie{n}\ln\sqrt{\gamma}$, where
$\gamma = \det(\gamma_{ij})$.  Positive $\meanextrinsic$ means the
spatial volume is contracting as the slicing moves forward; negative
$\meanextrinsic$ indicates expansion.  Taking the trace of
\cref{eq:adm-metric-evolution} (using the Jacobi formula
$\gamma^{ij}\,\partial_t\gamma_{ij} = \partial_t\ln\gamma$ and
metric compatibility $D_i\gamma^{jk} = 0$) gives the volume evolution
%
\begin{equation}
  \partial_t \ln\sqrt{\gamma}
  = -\lapse\,\meanextrinsic + D_i\beta^i \,,
  \label{eq:adm-volume-evolution}
\end{equation}
%
separating the physical contraction (from $\meanextrinsic$) and the
coordinate compression (from the divergence of the shift).  The mean
curvature appears prominently in the Hamiltonian constraint of
\cref{sec:adm-equations} and is the key scalar controlled by
gauge conditions such as maximal slicing ($\meanextrinsic = 0$).
\physnote{A maximal slice ($\meanextrinsic = 0$) has stationary spatial
volume.  Maximal slicing avoids the singularity-crashing behaviour of
geodesic slicing because the collapsing geometry produces large
$\meanextrinsic$, which the gauge condition then counteracts by
increasing the lapse (\cref{sec:bssn-gauge}).}

\paragraph{Codebase decomposition.}

The \texttt{ADMSpacetime} typeclass introduced in
\cref{sec:adm-lapse-shift} includes extrinsic curvature alongside the
three kinematic fields.  The spatial Christoffel symbols are also
provided, since the evolution equations require spatial covariant
derivatives:
%
\begin{lstlisting}[language=Haskell]
  extrinsicCurvature :: s -> Point s -> Sym3 'Covariant
  spatialChristoffel :: s -> Point s -> Christoffel3
\end{lstlisting}
%
The return type \texttt{Sym3 'Covariant} enforces at the type level that
$\extrinsic$ is a symmetric, purely covariant three-tensor, consistent
with the properties derived above.  For an analytic spacetime the method
returns a closed-form expression; for a numerical evolution it returns
current grid values.
